
	\section{Introduction}
	% \IEEEPARstart{T}{his} file~\cite{Baruah2015MultiprocessorSF}.


	% Introduce big problem
	% •	Software supports robotics
	% •	Real-world robot systems need to run multiple tasks simultaneously
	% •	In embedded systems, computation resource is limited, but performance guarantees often arise in many safety-critical applications/real-time systems
	% •	Therefore, scheduling and resource allocation algorithms are the key for such guarantee

	% •	Robotics:
	% o	Violating DDL is a new safety issue when computation resources are limited, and providing such support in the robotic community is less studied;
	% o	A good system should have both safety and performance, and providing such guarantee is important in safety-critical systems:
	% Hard guarantee has been studied in hard RTS;
	% Soft guarantee is more flexible, general, and allows better resource utilization
	% o	No prior work has studied safety and performance guarantee under dynamic environment context
	% o	Side note: “real-time” support/features are often about running fast rather than providing system-level performance and safety guarantees;

	% Dynamic environments
	% •	Robot is operated under changing/dynamic environment
	% •	Computation tasks’ input highly depends on environment, and its internal states
	% •	Resource requirements vary both temporally and spatially
	% •	To guarantee safety and performance, scheduling algorithms should be optimized based on resource needs
	% Available/related work/knowledge gap


	% In this work, we propose new scheduling/resource allocation algorithms to provide such guarantee under dynamic environments:
	% •	Safety and performance analysis, new SP metric based on previous work
	% •	SP optimization: priority assignments and ET optimization
	% Experiments show good results, summarize contributions.



	As robotics systems continue to evolve and proliferate, they increasingly rely on software and
	embedded systems to perform complex and mission-critical operations. Real-world robotic
	applications, from autonomous vehicles to industrial robots, often need to run multiple tasks
	simultaneously, each with specific performance and safety
	requirements~\cite{Heintzman2021-mw, Percept21_RTSS, Ash23RAL}
	% \rkwprev{Need to add several citations here; before submissions I will add several from my lab's
	% applications in aerial robots as they fit this work very well.  You should look for citations in
	% other applications that you list}. 
	However, embedded systems—the backbone of these robotic
	platforms—typically operate under constrained computational resources. Many safety-critical
	applications require not only high performance but also stringent performance and safety
	guarantees. Achieving such guarantees necessitates a careful and comprehensive modeling of
	computational systems.

    A significant safety issue arises in resource-constrained systems: it would be dangerous if 
	safety-critical tasks
    produce reliable outputs but take prohibitively long to do so. 
	For instance, consider an autonomous vehicle whose
    pedestrian detection algorithm generates a positive result only after the vehicle has already
    collided with a pedestrian. Addressing these challenges traditionally falls within the domain of
    real-time systems rather than robotics. However, computational resources influence not only
    low-level scheduling but also high-level robotic behaviors and performance
    \footnote{In this paper, performance of robotic tasks mostly refer to their output quality, such as more
    accurate localization for SLAM or shorter moving distance in terms of path planning}.
    Tasks' performance may be limited by the available computation resources.
    For example, anytime algorithms for high-level robot planning (\eg \cite{Huang23ICRA_TSP}) can
    achieve better performance if allocated additional computational resources (\ie longer runtime).
    However, the system will become unsafe if doing so causes safety-critical tasks to not have
    enough computation resources. This highlights the tight coupling between performance and safety,
    particularly when resources are limited. This paper advocates for an integrated approach to
    model these tasks and address these issues holistically and provide unified performance guarantees.


    The challenges grow even more complex when robots operate in dynamic
    environments. Under varying spatial and temporal conditions, tasks may follow different execution
    paths and require different amounts of computational
    resources~\cite{Ash23RAL, Sifat2021TowardsCA}.
    For example, SLAM may require a longer execution time when working under a crowded city
    environment than in a sparse rural environment.
    A simplistic but overly conservative solution might involve provisioning sufficient computational
    resources to handle the worst-case scenario. However, such an approach often results in resource
    waste and can negatively impact performance metrics such as battery
    life~\cite{Maxim2022ProbabilisticA}. Providing consistent safety and performance guarantees
    under stringent resource constraints and dynamic environments remains a critical yet
    underexplored research topic.

    Therefore, a good scheduling algorithm is necessary for the robotic systems to achieve good
    safety and performance. Task scheduling typically involves deciding the tasks' running order and
    computation resource allocation.
    In this paper, we propose a novel framework to analyze and optimize the safety and performance of
    robotic systems in dynamic environments. First, we enhance the Safety-Performance (SP) metric
    introduced by Sifat~\cite{Ash23RAL}, improving its formal definition and interpretability. Next,
    using the system-level SP metric as the objective function, 
	we propose online optimization algorithms to optimize task priority assignments and QoS
    configurations to dynamically allocate resources efficiently.
    To minimize the overhead of these online algorithms, we introduce an incremental optimization
    strategy that balances computational cost and performance. 
	The incremental optimization framework unifies the two major types of tasks, enviornment-dependent tasks 
	whose execution time distribution is decided by the external environment, and QoS-configurable tasks whose execution time
	are decided by the internal schedulers, then propose efficient strategies to perform optimization based on
	insights from real-time scheduling and numerical optimization theorems.

    We evaluated our framework in both real-world and simulation experiments. 
	The experiments demonstrated significant performance improvements against baseline methods,
	close-to-optimal solution quality, and very lightweight overhead.


	\textbf{Major contributions} are summarized as follows:
	\begin{itemize}
		\item We enhance the SP metric proposed in \cite{Ash23RAL}, improving its formal definition
        and interpretability.
        \item We formulate an online optimization problem to improve the SP metric under dynamic
        environmental conditions for systems with environment-dependent tasks and QoS-configurable tasks.
        \item We propose novel incremental optimization algorithms for task priority assignments and
        runtime configurations, achieving low-overhead performance improvements.
        \item We designed, implemented, and tested a robotic computing system on an embedded system,
        and released the code\footnote{\githubLink{}} to encourage further experiments.
	\end{itemize}


